\begin{description}
\item[(3.1)] Autumn (fall). The Moon, Earth and Sun are located on the same
  line; the Sun is seen from Earth in the constellation Virgo, next to the
  autumnal equinox point. \hfill(2 points)
\item[(3.2)] $598.8$ ($\approx 599$) million km. Juno, the Sun, the Earth
  and the Moon are nearly on the same line, so the distance is
  $3\,\mathrm{au} + 1\,\mathrm{au} + d(\mbox{Earth--Moon}) \approx
  4 \times 149.6 + 0.4 = 598.8$ million km.
  (2 points for $\pm 1$ million km, 1 point for $\pm 2$ million km.)
  \hfill(maximum 2 points)
\item[(3.3)] Sunset: about 3--4 days later (in mean solar days). The length
  of the day on the Moon equals the synodic month, 29.53 days (708.72
  hours), so the Sun moves $360^\circ/29.53 = 12.2^\circ$ per day on the
  lunar sky from east to west. The Sun is $40^\circ$ above the horizon and
  the site is on the lunar equator, so $40^\circ/12.2^\circ$ days
  $= 3.28$ days $= 78.7$ hours are needed for setting. (The extent of the
  solar disc is not considered; that adds about half an hour to complete
  sunset.) (Between 3--4 days: 4 points; 2 or 5 days: 2 points.)
  \hfill(maximum 4 points)\\
  Earth set: never. The Earth--Moon system is in synchronous rotation; the
  observing site is on the near side of the Moon, outside the libration
  zone, so the Earth is continuously observable. \hfill(2 points)
\item[(3.4)] Latitude $0^\circ$, longitude $50^\circ$\,E. If we sat at the
  origin of the selenographic coordinate system, the solar eclipse by the
  Earth would happen at the zenith; since the event is seen almost exactly
  towards the west at a horizontal altitude of $40^\circ$, the site is on
  the lunar equator, about $50^\circ$ east of the origin.
  (Latitude within $\pm 2^\circ$ = 3 points, $\pm 5^\circ$ = 2 points;
  longitude within $\pm 2^\circ$ ($48^\circ$--$52^\circ$) = 3 points,
  $\pm 5^\circ$ = 2 points.) \hfill(maximum 6 points)\\
  The area is Mare Fecunditatis (``Sea of Fecundity'' or ``Sea of
  Fertility''). \hfill(6 points)
\item[(3.5)] About 800 km. Disregarding the latitudinal difference of the
  Apollo-11 landing site from the observing site (it is less than
  $1^\circ$), both sites can be taken on the lunar equator, a great circle
  of circumference $D_{\mathrm{Moon}}\pi = 3474\,\mathrm{km} \times 3.14 =
  10914$ km. This corresponds to $360^\circ$, and the longitudinal
  difference between the sites is $50^\circ - 23.473^\circ = 26.527^\circ$,
  so the arc length is 803.7 km.
  ($\pm 100$ km: 6 points; $\pm 200$ km: 5 points; $\pm 300$ km: 4 points;
  $\pm 400$ km: 3 points.) \hfill(maximum 6 points)
\end{description}
